\section{Information conditions and controls}
\label{sec:information}

The three model conditions share the model, atomic proposal schema, editable surface, and evaluation protocol. The scalar condition receives the current task, harness, and scalar probe outcomes. The memory condition adds cumulative historical probe memory, including prior failures and behavioral differences that can inform a later decision. The diagnostic condition adds a compact execution trace to that memory, exposing failure categories, mechanism differences, and cost information. Select outcomes remain outside every proposer input.

Two non-model references complete the design. Random legal atomic edits test how much of the observed change can be obtained by sampling the same proposal space with an independent random seed and no language-model calls. The fixed harness measures the unchanged starting process under the same final evaluator. The memory-minus-random contrast is consequently an end-to-end comparison of research processes, whereas memory-minus-scalar and diagnostic-minus-memory isolate incremental information under their respective parent conditions.

Earlier panels played distinct roles in the study design. V001 tested diagnostic continuation and obtained a median paired difference of 0.0, yielding an exploratory development-panel null. V002 tested cumulative memory and obtained a median difference of 0.001224, also yielding an exploratory development-panel null. V003 preserved those records and introduced explicit atomic semantics, a random legal-edit reference, four continuation decisions per arm, six fresh tasks, and independent model-call replicates. These design changes motivate the present panel without converting the earlier exploratory outcomes into confirmatory support.
